\subsection{1. Literaturrecherche:} Identifikation von State-of-the-Art Ansätzen für synthetische Daten und Sensor-Schnittstellen.
\begin{enumerate}[label=1.\arabic*]
    \item Welche Technologien werden verwendet, um strukturierte klinische Daten synthetisch zu erzeugen?
    \item Welche Technologien werden verwendet um Sensor-Daten in das KIS einzubinden?
    \item Welche Anforderungen werden in der Literatur an synthetische strukturierte klinische Daten gestellt?
\end{enumerate}

\subsection{2. Definition Qualitätskriterien:} Festlegung didaktischer und technischer Anforderungen für das Heidelberger Lehre-KIS.
\begin{enumerate}[label=2.\arabic*]
    \item Befragung von Expert*innen
    \item Analyse der Ergebnisse: Welche Anforderungen haben Spezialisten an die Daten?
\end{enumerate}

\subsection{3. Prototypische Implementierung:} Erzeugung von Datensätzen und Anbindung von Wearables/Hardware und Einbindung in das KIS.
\begin{enumerate}[label=3.\arabic*]
    \item Wie könnten strukturierte klinische Daten synthetisch durch nicht-Sprachmodell-basierte Agenten erzeugt werden?
    \item Wie könnten strukturierte klinische Daten synthetisch durch Sprachmodelle erzeugt werden?
    \item Wie können Datensätze (sowohl sensorische und synthetische) im KIS eingebunden werden?
    \item Wie können die Daten dynamisch ("lebendig") in ein KIS eingebunden werden?
\end{enumerate}

\subsection{4. Bewertung:} Evaluierung der Ansätze hinsichtlich Eignung, Komplexität und Ressourcenverbrauch.
\begin{enumerate}[label=4.\arabic*]
    \item Aus welchen Quellen sollen die Patientendaten für die Evaluierung stammen? 
    \item Wie können die verschiedenen Patientendaten einheitlich dargestellt werden?
    \item Wie kann die Evaluierung effizient durchgeführt werden?
    \item Gibt es statistische Unterschiede zwischen den Patientengruppen?
    \item Welche qualitativen Auffälligkeiten zeigen sich zwischen den Patientengruppen?
\end{enumerate}
